Statevector simulation of quantum circuits is limited by memory bandwidth, and data-centre FPGAs with high-bandwidth memory offer bandwidth comparable to workstation GPUs. We present a statevector simulator for the AMD Alveo U55C that executes every gate with one multiplexed $2\times2$ datapath primitive applied to on-chip tiles of the state, runs one streaming engine per HBM pseudo-channel and performs qubit remaps on the card. A single card simulates 30 qubits in single precision. We measure the memory ceilings of the card (\NumHbmThirtytwoCopy~GB/s over 32 pseudo-channels, \NumHbmSixteenPerCURead~GB/s per channel), study engines, channels, port width, tiling, gate fusion and exchange, and verify correctness with a fidelity suite of \NumSuiteTotal{} cases started from entangled states, all of which pass. Compared with an existing HBM kernel for the same card, end-to-end time falls by \NumSpeedupVsevenMin to \NumSpeedupVsevenMax$\times$ at 24 qubits and energy per circuit by \NumEnergyVsevenRatioMin to \NumEnergyVsevenRatioMax$\times$. Against Qiskit Aer, qsim and Qulacs on a 128-core server, the card is faster for bandwidth-bound circuits from \NumXplatGhzCrossN{} qubits (GHZ-30 in \NumFlagGhzThirtyEtoE~s, \NumXplatRatioGhzThirty$\times$ faster than the best CPU configuration), while for compute-bound circuits (QFT, random circuits, VQE ansatz) the server remains faster, needing \NumXplatRatioOtherThirtyMin to \NumXplatRatioOtherThirtyMax{} of the FPGA time at 30 qubits. The card draws \NumFlagPowerMin to \NumFlagPowerMax~W while simulating. An NVIDIA A100 with Qiskit Aer and CUDA-Q is faster than the card for compute-bound circuits (by \NumXplatBestGpuRatioRqcThirty~to \NumXplatBestGpuRatioQftThirty$\times$ for random circuits and QFT at 30 qubits), about as fast for GHZ, and holds up to \NumGpuCudaqSingleMaxN~qubits.
